% !TeX root = ../main.tex
% ============================================================
%  algorithms/alg_bounding_procedure.tex  —  Main DBP
% ============================================================

\begin{algorithm}[H]
\DontPrintSemicolon
\caption{Deterministic Bounding Procedure (\DBP)}
\label{alg:dbp}

\KwIn{Tax bounds $[\taxindexlb,\taxindexub]$; tolerance $\epsilon>0$.}
\KwOut{Tax index $\taxindex^{\dagger}$ and bilevel-feasible response
  $\bz^{\dagger}$ with $\gap\le\epsilon$.}

\BlankLine
\tcp{Initialisation. The relaxation master optimises over the same
  feasible region $\setX$ as the follower (\Cref{sec:model:follower}),
  which does \emph{not} depend on $\taxindex$ -- only the follower's
  objective~\eqref{eq:follower} does. This is what lets a cut generated
  at one tax index remain valid at every other (\Cref{lem:lb_validity}).}
$\LB\leftarrow-\infty$; $\UB\leftarrow+\infty$; $\mathcal{C}\leftarrow\emptyset$\;
$\taxindex^{(0)}\leftarrow$ Stage~1/Stage~2 selection of
  \Cref{sec:method:tiebreak} applied to $\mathcal{C}=\emptyset$: the
  smallest tax index attaining
  $\min_{\taxindex\in\mathcal{T}^{\tau},\,\bz\in\setX}\Emis(\bz)$\;
$k\leftarrow 0$\;

\BlankLine
\While{$\UB-\LB > \epsilon\cdot\UB$}{%
  \tcp{---- Upper bound: exact follower evaluation at $\taxindex^{(k)}$ ----}
  % !TeX root = ../main.tex
% ============================================================
%  algorithms/alg_upper_bound.tex  —  UB subroutine (inlined)
% ============================================================

% NOTE: this file is \input inside alg_bounding_procedure.tex;
%       it does NOT open a new algorithm float.

\tcp*[h]{Evaluate the follower exactly at the master's current candidate $\taxhat^{(k)}$}
Solve~\eqref{eq:follower} at $\taxhat=\taxhat^{(k)}$ to optimality, and
among its optimal responses select the one with least emissions
(\Cref{as:opt}), giving $\bz^{(k)}\in\Sol{\taxhat^{(k)}}$,
$c_{k}\defeq\Cost\!\bigl(\bz^{(k)}\bigr)$,
$e_{k}\defeq\Emis\!\bigl(\bz^{(k)}\bigr)$\;
\If{$e_{k}<\UB$}{%
  $\UB\leftarrow e_{k}$;\quad
  $\taxhat^{\dagger}\leftarrow\taxhat^{(k)}$;\quad
  $\bz^{\dagger}\leftarrow\bz^{(k)}$\;
}
\tcp{Record the cut
  $\Cost(\bz)+\tfrac{\taxhat}{100}\Emis(\bz)\le c_{k}+\tfrac{\taxhat}{100}e_{k}$,
  valid for every $\taxhat\in\mathcal{T}^{\tau}$ since
  $\bz^{(k)}\in\setX$ regardless of $\taxhat$ (\Cref{lem:lb_validity})}
$\mathcal{C}\leftarrow\mathcal{C}\cup\bigl\{(\taxhat^{(k)},\,c_{k},\,e_{k})\bigr\}$\;
\;
  \BlankLine
  \tcp{---- Lower bound: re-solve the relaxation master with the new cut ----}
  % !TeX root = ../main.tex
% ============================================================
%  algorithms/alg_lower_bound.tex  —  LB subroutine (inlined)
% ============================================================

% NOTE: this file is \input inside alg_bounding_procedure.tex;
%       it does NOT open a new algorithm float.

\tcp*[h]{Stage 1: re-solve the relaxation master's emissions minimum with all cuts collected so far}
$\LB\leftarrow
  \displaystyle\min_{\taxindex\in\mathcal{T}^{\tau},\,\bz\in\setX}
  \Bigl\{\Emis(\bz)\ :\
  \Cost(\bz)+\tfrac{\taxindex}{100}\Emis(\bz)\le c_{j}+\tfrac{\taxindex}{100}e_{j}
  \ \ \forall(\taxindex^{(j)},c_{j},e_{j})\in\mathcal{C}\Bigr\}$\;
\tcp*[h]{Stage 2: among master points attaining $\LB$, take the smallest tax index (\Cref{lem:tiebreak_validity})}
$\taxindex^{(k+1)}\leftarrow
  \displaystyle\min\Bigl\{\taxindex\in\mathcal{T}^{\tau}\ :\ \exists\,\bz\in\setX,\ \Emis(\bz)=\LB,\
  \Cost(\bz)+\tfrac{\taxindex}{100}\Emis(\bz)\le c_{j}+\tfrac{\taxindex}{100}e_{j}
  \ \ \forall(\taxindex^{(j)},c_{j},e_{j})\in\mathcal{C}\Bigr\}$\;
\;
  \BlankLine
  $\gap\leftarrow (\UB-\LB)/\UB$\;
  $k\leftarrow k+1$\;
}
\BlankLine
\Return $\taxindex^{\dagger},\,\bz^{\dagger},\,\gap$\;
\end{algorithm}
